\clearpage
\VerdierPhysicalPageAnchor{0240}%
\begingroup\small\setlength{\parskip}{0pt}\setlength{\emergencystretch}{1.5em}
\setlength{\abovedisplayskip}{4pt plus 2pt minus 2pt}%
\setlength{\belowdisplayskip}{4pt plus 2pt minus 2pt}%

\VerdierSemanticAnchor{ch4d-4-4-5}
\noindent\textbf{4.4.5.}\quad
The data are the same as in (4.4.4), but in addition an object $X$ of
$D(\mathcal A)$ is given. Let $\bigl(X(p,q),\delta(p,q,r)\bigr)$ be the
second spectral object with values in $D(\mathcal A)$ associated with the
complex $X$ (4.3.1). The functor $F$ takes this spectral object to a
spectral object with values in $\mathcal B$. One thereby obtains a spectral
sequence functorial in $F$ and in the argument $X$ as the latter varies in
the category $D(\mathcal A)$ (4.3.3). This spectral sequence is called the
\emph{second spectral sequence of the functor $F$ relative to the complex
$X$}. For this spectral sequence we always use the indexing introduced in
(Chapter~II, 4.3.3.2).

\VerdierSemanticAnchor{ch4d-4-4-6}
\medskip\noindent\textbf{Proposition 4.4.6.}\quad

\VerdierSemanticAnchor{ch4d-4-4-6-a}
\noindent\textit{a) The term $\mathrm{II}_2^{p,q}$ of the second spectral
sequence of $F$ relative to the complex $X$ is given by}
\[
\mathrm{II}_2^{p,q}=F^p\bigl(H^q(X)\bigr).
\]

\VerdierSemanticAnchor{ch4d-4-4-6-b}
\noindent\textit{b) The terms $\mathrm{II}_{\infty}^{p,q}$ are the
associated graded objects of the suitably filtered objects $F^{p+q}(X)$.}

\VerdierSemanticAnchor{ch4d-4-4-6-proof}
Indeed, by (4.3.4.1), there is a canonical isomorphism
\[
X(p,p+1)\xrightarrow{\sim}H^p(X)[-p],
\]
whence, using (Chapter~II, 4.3.3.2),
\[
\mathrm{II}_2^{p,q}
=F^{p+q}\bigl(H^q(X)[-q]\bigr)
=F^p\bigl(H^q(X)\bigr),
\]
which gives the (abusive) equality in a). Assertion b) follows from
(Chapter~II, 4.3.2).

\VerdierSemanticAnchor{ch4d-4-5}
\subsection{Convergence problems.}

\VerdierSemanticAnchor{ch4d-4-5-1}
\noindent\textbf{Definition 4.5.1.}\quad
Let $\mathcal A$ and $\mathcal B$ be two abelian categories. A
cohomological functor $F\colon D^*(\mathcal A)\longrightarrow\mathcal B$
($*=+,-,b,{}$ ``empty'') is said to be \emph{stationary on the right}
(respectively, \emph{on the left}) if there is an integer $n_0$ such that,
for every object $X$ of $D^*(\mathcal A)$ satisfying $H^i(X)=0$ for every
$i<0$ (respectively, $>0$), one has $F^n(X)=0$ for every $n<n_0$
(respectively, $n>n_0$). The functor $F$ is said to be \emph{stationary} if
it is stationary on the right and on the left. An exact functor
$G\colon D^*(\mathcal A)\longrightarrow D^{*'}(\mathcal B)$ is said to be
\emph{stationary on the right} (respectively, \emph{stationary on the
left}, respectively, \emph{stationary}) if the cohomological functor
obtained by composing $G$ with the canonical cohomological functor
$D^{*'}(\mathcal B)\longrightarrow\mathcal B$ is stationary on the right
(respectively, stationary on the left, respectively, stationary).
\endgroup
